\subsubsection{Mass and reference densities}
\label{sec:invariant-mass-reference-densities}

In dimension three, an invariant positive background has a nonnegative complex Hessian. This determines its upper Monge--Amp\`ere mass and the top mass of its singular coefficient limits. On the Iwasawa quotient, the same calculation identifies exactly which smooth reference probabilities retain the determinant formula.

Retain the invariant coframe, $h_0$, $H$ and $\theta_A$ from \cref{sec:parallelizable-gauduchon-cones}, now with $n=3$. Write $\mu_0=h_0^3/H$ to distinguish this invariant probability from other reference measures, and let $\mathscr M$ be the smooth positive probabilities on $X$.

\Needspace{12\baselineskip}
\begin{proposition}\label{prop:invariant-threefold-upper-mass}\index{Monge--Ampere mass@Monge--Amp\`ere mass!invariant threefold backgrounds}
For $A>0$ on any such compact complex Lie-group quotient,
\begin{equation}\label{eq:invariant-threefold-upper-mass}
 v_+(\theta_A)=H\det A
      =\max_{\nu\in\mathscr M}\mathcal V_\nu(\theta_A).
\end{equation}
The invariant probability $\mu_0$ attains this maximum. In particular the quotient has the bounded mass property.
\end{proposition}

\begin{proof}
The coframe has holomorphic structure coefficients, so
\begin{equation}\label{eq:invariant-pluripositive-form}
 \ddc\theta_A=\frac1{4\pi^2}
       \sum_{j,k}A_{jk}\,\mathrm d\alpha_j\wedge\mathrm d\bar\alpha_k\geq0.
\end{equation}
Indeed $\bar\partial(\alpha_j\wedge\bar\alpha_k)
=-\alpha_j\wedge\mathrm d\bar\alpha_k$, giving the displayed sign after applying $\partial$ and multiplying by $(\mathrm i/(2\pi))^2$. Diagonalizing the positive coefficient matrix expresses the result as a sum of nonnegative multiples of $\beta\wedge\bar\beta$ for invariant $(2,0)$-forms $\beta$. In dimension three these are nonnegative $(2,2)$-forms: their wedge with $\mathrm i\xi\wedge\bar\xi$ has the squared modulus of $\xi\wedge\beta$ as coefficient.

By \cref{prop:parallelizable-gaud-volume}, $\theta_A$ is balanced, hence Gauduchon. The maximum case of \cref{thm:threefold-mass-minimizer} therefore gives $v_+(\theta_A)=\int_X\theta_A^3=H\det A$. The density supremum equals this upper mass by \cref{thm:gaud-volume-mass-extrema}, and \eqref{eq:parallelizable-gaud-volume} gives attainment at $\mu_0$. Finally \cref{prop:bmp-metric} makes the finite upper mass a property of the manifold.
\end{proof}

\Needspace{12\baselineskip}
\begin{corollary}\label{cor:invariant-singular-coefficient-zero-mass}\index{Monge--Ampere mass@Monge--Amp\`ere mass!singular invariant coefficients}
If $A\geq0$ is singular, then every bounded $u\in\PSH(X,\theta_A)$ has
\[
 \int_X(\theta_A+\ddc u)^3=0.
\]
For every possibly unbounded $u\in\PSH(X,\theta_A)$, its non-pluripolar top product is also zero. Thus the nonzero semipositive coefficient boundary currents from \cref{prop:parallelizable-gaud-volume} have zero non-pluripolar top mass.
\end{corollary}

\begin{proof}
For $\varepsilon>0$, the same bounded $u$ is $\theta_{A+\varepsilon I}$-plurisubharmonic. The preceding proposition bounds its mass by $H\det(A+\varepsilon I)$, which tends to zero. By the multilinear bounded construction in \cref{prop:bounded-products},
\[
 (\theta_{A+\varepsilon I}+\ddc u)^3
 =\bigl(\theta_A+\ddc u+\varepsilon h_0\bigr)^3
\]
is a polynomial in $\varepsilon$ with finite mixed-current coefficients. Its total mass converges to the asserted mass. Positivity then makes the limiting measure zero.

For unbounded $u$, subtract a constant so that $u\leq0$. Since $\theta_A\geq0$, the canonical cutoffs $u_k=\max(u,-k)$ are bounded and $\theta_A$-plurisubharmonic. Their top products vanish by the bounded assertion. Restricting them to $\{u>-k\}$ and taking the defining non-pluripolar limit gives zero as well.
\end{proof}

Now specialize to the Iwasawa quotient of \cref{ex:iwasawa-complete-gaud-cone}. Put
\[
 E=\mathbb C/\mathbb Z[\mathrm i],\qquad B=E^2,\qquad
 p:X\longrightarrow B,\quad [z]\longmapsto([z_1],[z_2]).
\]
This is a holomorphic submersion with connected elliptic fibers. The coordinate description of the fundamental cube gives $p_*\mu_0=\mu_B$, where $\mu_B$ is the Haar probability on $B$. A smooth probability is \emph{fiber invariant} if it has the form
\[
 \mu_f=(f\circ p)\mu_0,\qquad f\in C^\infty(B),\quad
 f>0,\quad\int_B f\,\mu_B=1.
\]
Only invariance along the central elliptic fibers is required here.

\Needspace{12\baselineskip}
\begin{theorem}\label{thm:iwasawa-fiber-invariant-density}\index{Gauduchon variational volume!fiber-invariant densities}
For every $A>0$ and fiber-invariant $\mu_f$,
\begin{equation}\label{eq:iwasawa-fiber-invariant-volume}
 \mathcal V_{\mu_f}(\theta_A)=H\det A,
 \qquad H=\frac6{\pi^3}.
\end{equation}
Its unique minimizing Gauduchon metric normalized by $D_{\mu_f}=1$ is
\begin{equation}\label{eq:iwasawa-fiber-invariant-minimizer}
 \gamma_{A,f}=\frac{\theta_A+\ddc(p^*v_{A,f})}{(H\det A)^{1/3}},
 \qquad\gamma_{A,f}^3=\mu_f,
\end{equation}
where $v_{A,f}$ solves a closed positive Monge--Amp\`ere equation on $B$ specified in the proof. Conversely, for a fixed $A>0$, the smooth positive probabilities maximizing $\mathcal V_\nu(\theta_A)$ are exactly the fiber-invariant ones.

For each fixed real linear $W\subset\operatorname{Herm}(3)$ containing a positive matrix, the Hessian metric for $\mu_f$ is
\[
 g_A(U,Z)=\frac13\operatorname{tr}(A^{-1}UA^{-1}Z).
\]
Its unit-volume hypersurface $S_W=\{A\in W:A>0,\ H\det A=1\}$ is complete and proper. The coefficient Cauchy criterion and intrinsic bounds of \cref{thm:parallelizable-unit-volume-completion} hold with $n=3$.
\end{theorem}

\begin{proof}
The map $p$ is well-defined because lattice multiplication adds Gaussian integers to the first two coordinates. It is a submersion, and fixing representatives of those coordinates identifies each fiber with the central quotient $E$. Integrating first in the third coordinate of the fundamental cube proves the pushforward identity.

\emph{The closed equation on the base.}
Write the positive Hermitian matrix in block form and take its Schur complement:
\[
 \begin{gathered}
 A=\begin{pmatrix}C&b\\ b^*&a\end{pmatrix},\\
 a=A_{33}>0,\qquad S=C-a^{-1}bb^*>0.
 \end{gathered}
\]
On $B$, let
\[
 \omega_S=\frac{\mathrm i}{2\pi}\sum_{j,k=1}^2
       S_{jk}\,\mathrm dz_j\wedge\mathrm d\bar z_k,
 \qquad Q_B=\int_B\omega_S^2=\frac2{\pi^2}\det S.
\]
Apply \cref{thm:tw-smooth} on this closed background with density $Q_Bf\mu_B$. Closed Stokes fixes the scalar to zero. It gives a smooth real $v=v_{A,f}$, unique up to constant, such that
\begin{equation}\label{eq:iwasawa-schur-base-equation}
 (\omega_S+\ddc v)^2=Q_Bf\mu_B,
 \qquad\omega_S+\ddc v>0.
\end{equation}
If $K$ is the coefficient matrix of $\ddc v$ in the base coframe, this says $\det(S+K)=(\det S)f$. The form
\[
 \Omega=\theta_A+\ddc(p^*v)
\]
has coefficient matrix $\left(\begin{smallmatrix}C+K&b\\b^*&a\end{smallmatrix}\right)$ in the invariant coframe. Schur positivity gives $\Omega>0$; its determinant is $a\det(S+K)=(\det A)(f\circ p)$. Consequently $\Omega^3=(H\det A)\mu_f$.

This solution is also balanced. Put $\beta=\ddc(p^*v)$, which is closed and uses only $\alpha_1,\alpha_2$ and their conjugates. Since $\mathrm d\alpha_3=-\alpha_1\wedge\alpha_2$, every term of $\partial\theta_A$ contains $\alpha_1\wedge\alpha_2$, and every conjugate term contains $\bar\alpha_1\wedge\bar\alpha_2$. Thus $\mathrm d\theta_A\wedge\beta=0$. Together with $\mathrm d\theta_A^2=0$, this proves
\[
 \mathrm d\Omega^2=\mathrm d\theta_A^2+2\mathrm d\theta_A\wedge\beta=0.
\]
The smooth dual formula \cref{thm:gaud-volume-dual-formula} now gives \eqref{eq:iwasawa-fiber-invariant-volume}. Its equality characterization, the balancedness of $\Omega$, and degree-two homogeneity of $D_{\mu_f}$ give the unique metric \eqref{eq:iwasawa-fiber-invariant-minimizer}.

\emph{All maximizing densities.}
Let $Z_3=\partial/\partial z_3$ be dual to $\alpha_3$, and set $\tau_j=(\mathrm i/(2\pi))\alpha_j\wedge\bar\alpha_j$. Directly from \eqref{eq:invariant-pluripositive-form},
\[
 \ddc\theta_A=A_{33}\tau_1\wedge\tau_2,
 \qquad h_0^3=6\tau_1\wedge\tau_2\wedge\tau_3.
\]
For every smooth real $u$, \eqref{eq:threefold-mass-dirichlet} therefore becomes
\begin{equation}\label{eq:iwasawa-vertical-mass-defect}
 \int_X(\theta_A+\ddc u)^3
   =H\det A-\frac{A_{33}H}{2}\int_X|Z_3u|^2\mu_0.
\end{equation}
Here $\mathrm du\wedge\dc u=(\mathrm i/(2\pi))\partial u\wedge\bar\partial u$; wedging with $\tau_1\wedge\tau_2$ selects its third coefficient. This also checks the factor $1/2$.

For an arbitrary smooth positive probability $\nu$, let $u_\nu$ be its smooth normalized solution. Its mass equals $\mathcal V_\nu(\theta_A)$ by the dual formula. Equality with $H\det A$ in \eqref{eq:iwasawa-vertical-mass-defect} holds precisely when $Z_3u_\nu=0$. As $u_\nu$ is real, both real directions of each elliptic fiber then vanish. Connectedness of the fibers gives $u_\nu=p^*v$ for a smooth $v$ on $B$. Its determinant relative to $\mu_0$ is a function on $B$, so the normalized equation makes $\nu$ fiber invariant. The converse was proved by the base construction.

\emph{Completion in the background parameters.}
For fixed $f$, formula \eqref{eq:iwasawa-fiber-invariant-volume} holds on the entire positive matrix cone. Differentiating its logarithmic root gives the stated Hessian metric. The matrix-path and coefficient-limit proof of \cref{thm:parallelizable-unit-volume-completion} applies unchanged to this identical metric on each fixed $W$. It proves the intrinsic bounds, completeness, properness and Cauchy criterion, without requiring the minimizing forms to have constant coefficients.
\end{proof}

There is a useful distinction in the minimizing-metric estimates. For a form $\eta$ with coefficient matrix $B(x)$ in the invariant coframe, write $\overline\eta=\int_XB(x)\mu_0$ for its coefficient mean. The mean-zero derivative identity \eqref{eq:parallelizable-zero-mean} gives
\[
 \overline{\gamma}_{A,f}=\frac{A}{(H\det A)^{1/3}}.
\]
Thus on $S_W$ the spectral distance bounds control these means, and a normalized nonzero singular coefficient limit makes their largest eigenvalue diverge. The corresponding minimizing metrics cannot be uniformly bounded above by a fixed multiple of $h_0$. For $f\ne1$, this mean estimate does not assert the pointwise order bound \eqref{eq:parallelizable-minimizer-order} for the varying forms $\gamma_{A,f}$.

Reference densities varying along the fibers give strict loss of volume. Choose a smooth real $\chi$ supported in a coordinate ball with $Z_3\chi\not\equiv0$, and let $t\ne0$ be sufficiently small that $\theta_A+t\ddc\chi>0$. Put
\[
 M(t)=\int_X(\theta_A+t\ddc\chi)^3,
 \qquad\nu_t=\frac{(\theta_A+t\ddc\chi)^3}{M(t)}.
\]
Then the same potential solves the equation for $\nu_t$, and
\[
 \mathcal V_{\nu_t}(\theta_A)
 =H\det A-\frac{A_{33}H}{2}t^2\int_X|Z_3\chi|^2\mu_0
 <H\det A.
\]
The reference $\nu_t$ here is constructed for the fixed background $A$. This exhibits measure dependence; it does not determine the Hessian completion for an arbitrary fixed density varying along the fibers.
